\section*{Problemas}

\subsection*{Enunciados de los problemas}

% Recordar
\begin{problema}
	\label{prob:c0cebb4}
	El tiempo de ensamblaje de un componente mecánico sigue una distribución con media poblacional $\mu = 45$ minutos y desviación estándar $\sigma = 12$ minutos. Si se toma una muestra aleatoria de tamaño $n$, calcula el error estándar de la media muestral, denotado por $\text{DE}(\bar{X}) = \sigma_{\bar{X}}$, para cada uno de los siguientes tamaños de muestra y comenta sobre la velocidad de reducción del error:
	\begin{enumerate}
		\item $n = 9$
		\item $n = 36$
		\item $n = 144$
		\item $n = 576$
	\end{enumerate}
\end{problema}

% Comprender
\begin{problema}
	\label{prob:116017b}
	Una población finita consta de $N = 6$ elementos distintos $\{a, b, c, d, e, f\}$. Se desea extraer una muestra aleatoria de tamaño $n = 3$.
	\begin{enumerate}
		\item Determina el número total de muestras posibles si la extracción se realiza \textbf{con reemplazo}.
		\item Determina el número total de muestras posibles si la extracción se realiza \textbf{sin reemplazo} y sin importar el orden (muestras no ordenadas).
		\item Explica formalmente por qué, a medida que el tamaño poblacional $N$ crece hacia el infinito, la diferencia probabilística entre muestreo con y sin reemplazo tiende a cero.
	\end{enumerate}
\end{problema}

% Aplicar
\begin{problema}
	\label{prob:2da8cc3}
	Una embotelladora de agua mineral desea calibrar su línea de producción. La cantidad vertida en cada botella es una variable aleatoria con desviación estándar conocida de $\sigma = 8$ mililitros. ¿Cuál es el tamaño de muestra mínimo $n$ que un ingeniero de calidad debe inspeccionar para tener una certeza probabilística de al menos $95\%$ de que la media muestral $\bar{X}$ no difiera de la verdadera media poblacional de llenado $\mu$ en más de $1.5$ mililitros?
\end{problema}

% Analizar
\begin{problema}
	\label{prob:0c980d4}
	Demuestra algebraicamente por qué el estimador de la varianza muestral requiere dividir entre $n-1$ (y no entre $n$) para ser insesgado. Sea $S^2 = \frac{1}{n-1} \sum_{i=1}^n (X_i - \bar{X})^2$ la varianza muestral de una MAS obtenida de una población con media $\mu$ y varianza $\sigma^2$. Demuestra paso a paso que $E(S^2) = \sigma^2$.
	\emph{Sugerencia:} Suma y resta $\mu$ dentro del término cuadrático $(X_i - \bar{X})^2 = [(X_i - \mu) - (\bar{X} - \mu)]^2$.
\end{problema}

% Evaluar
\begin{problema}
	\label{prob:293fd20}
	El Teorema del Límite Central postula que la suma de variables independientes convergerá a la distribución normal siempre que el momento de segundo orden ($\sigma^2$) sea finito. Evalúa qué sucede cuando se extraen muestras de una población con colas extremadamente pesadas: la distribución de Cauchy estándar, cuya densidad es $f(x) = \frac{1}{\pi(1+x^2)}$ para $x \in (-\infty, \infty)$.
	\begin{enumerate}
		\item Demuestra o argumenta formalmente por qué la integral que define la esperanza teórica $E(X)$ no converge absolutamente en la distribución de Cauchy.
		\item Si se extrae una muestra aleatoria de tamaño $n$ de una población Cauchy, ¿cuál es la función de densidad o comportamiento de la media muestral $\bar{X}_n$? ¿Ocurre convergencia hacia la normal conforme $n \to \infty$? Juzga qué implicación práctica tiene esto para un analista que asuma ciegamente el TLC en cualquier conjunto de datos.
	\end{enumerate}
\end{problema}

% Crear
\begin{problema}
	\label{prob:b7567ec}
	Diseñe un problema original de aplicación del Teorema del Límite Central: elija una población con media y varianza poblacional conocidas (distinta del dado o la embotelladora ya vistos en esta sección), un tamaño de muestra $n \ge 30$, y calcule la probabilidad de que la media muestral se encuentre dentro de un rango específico que usted defina.
\end{problema}

\subsection*{Soluciones detalladas}

% Recordar
\begin{solproblema}[prob:c0cebb4]
	El error estándar de la media ($\sigma_{\bar{X}}$) mide la dispersión de los promedios muestrales alrededor de $\mu$ y está dado por $\sigma_{\bar{X}} = \frac{\sigma}{\sqrt{n}}$. Con $\sigma = 12$:
	\begin{enumerate}
		\item Para $n = 9$: $\sigma_{\bar{X}} = \frac{12}{\sqrt{9}} = 4.00$ min.
		\item Para $n = 36$: $\sigma_{\bar{X}} = \frac{12}{\sqrt{36}} = 2.00$ min.
		\item Para $n = 144$: $\sigma_{\bar{X}} = \frac{12}{\sqrt{144}} = 1.00$ min.
		\item Para $n = 576$: $\sigma_{\bar{X}} = \frac{12}{\sqrt{576}} = 0.50$ min.
	\end{enumerate}
	\textbf{Comentario sobre la velocidad de reducción:} la dispersión decrece de manera estrictamente inversamente proporcional a la raíz cuadrada del tamaño muestral ($\sqrt{n}$). Para reducir el error estándar a la mitad (por ejemplo, de 4.0 a 2.0, o de 2.0 a 1.0), es obligatorio cuadruplicar el tamaño de la muestra inspeccionada ($9 \to 36 \to 144$).
\end{solproblema}

% Comprender
\begin{solproblema}[prob:116017b]
	\begin{enumerate}
		\item \textbf{Muestreo con reemplazo:} en cada una de las $n = 3$ extracciones hay exactamente $N = 6$ opciones disponibles, dado que el elemento extraído se devuelve a la población antes de la siguiente selección. Por el principio multiplicativo del conteo:
		\begin{align}
			N^n = 6^3 = 216 \text{ muestras.}
		\end{align}
		\item \textbf{Muestreo sin reemplazo (no ordenado):} el número total de subconjuntos posibles es el coeficiente binomial:
		\begin{align}
			\binom{N}{n} = \binom{6}{3} = \frac{6!}{3!(6-3)!} = 20 \text{ muestras.}
		\end{align}
		\item \textbf{Equivalencia asintótica cuando $N \to \infty$:} al extraer una muestra fija $n$ sin reemplazo, la probabilidad de seleccionar un elemento específico en la primera extracción es $1/N$, en la segunda es $1/(N-1)$, y en la $k$-ésima es $1/(N-k+1)$. A medida que $N \to \infty$ con $n$ constante, las razones $\frac{N-k}{N} \to 1$. Por lo tanto, el agotamiento poblacional es despreciable y el muestreo sin reemplazo converge en distribución al muestreo con reemplazo (independiente).
	\end{enumerate}
\end{solproblema}

% Aplicar
\begin{solproblema}[prob:2da8cc3]
	Datos: $\sigma = 8$ ml, margen de error máximo tolerado $E = 1.5$ ml, nivel de confianza del $95\%$ ($\alpha = 0.05$, $z_{\alpha/2}=1.96$). El margen de error de estimación de la media es $E = z_{\alpha/2} \cdot \sigma/\sqrt{n}$. Despejando el tamaño de muestra mínimo:
	\begin{align}
		n = \left( \frac{z_{\alpha/2} \cdot \sigma}{E} \right)^2 = \left( \frac{1.96 \times 8}{1.5} \right)^2 \approx (10.4533)^2 \approx 109.27.
	\end{align}
	Dado que $n$ debe ser un número entero y se requiere garantizar al menos un $95\%$ de certeza, se redondea siempre al entero superior inmediato:
	\begin{align}
		n_{\text{mín}} = 110 \text{ botellas.}
	\end{align}
\end{solproblema}

% Analizar
\begin{solproblema}[prob:0c980d4]
	Queremos encontrar $E(S^2)$ donde $S^2 = \frac{1}{n-1} \sum_{i=1}^n (X_i - \bar{X})^2$. Sumamos y restamos la verdadera media poblacional $\mu$ dentro del paréntesis:
	\begin{align}
		\sum_{i=1}^n (X_i - \bar{X})^2 &= \sum_{i=1}^n [(X_i - \mu) - (\bar{X} - \mu)]^2 \\
		&= \sum_{i=1}^n (X_i - \mu)^2 - 2(\bar{X} - \mu) \sum_{i=1}^n (X_i - \mu) + n(\bar{X} - \mu)^2.
	\end{align}
	Como $\sum_{i=1}^n (X_i - \mu) = n\bar{X} - n\mu = n(\bar{X} - \mu)$, sustituyendo en el término central:
	\begin{align}
		\sum_{i=1}^n (X_i - \bar{X})^2 &= \sum_{i=1}^n (X_i - \mu)^2 - 2n(\bar{X} - \mu)^2 + n(\bar{X} - \mu)^2 = \sum_{i=1}^n (X_i - \mu)^2 - n(\bar{X} - \mu)^2.
	\end{align}
	Aplicando el operador esperanza a ambos lados, con $E[(X_i - \mu)^2] = \sigma^2$ y $E[(\bar{X} - \mu)^2] = \sigma^2/n$:
	\begin{align}
		E\left[ \sum_{i=1}^n (X_i - \bar{X})^2 \right] = n\sigma^2 - n\left(\frac{\sigma^2}{n}\right) = n\sigma^2 - \sigma^2 = (n-1)\sigma^2.
	\end{align}
	Finalmente, dividiendo por $(n-1)$:
	\begin{align}
		E(S^2) = \frac{(n-1)\sigma^2}{n-1} = \sigma^2.
	\end{align}
	\textbf{Conclusión:} si se dividiera entre $n$, el estimador tendría una esperanza de $\frac{n-1}{n}\sigma^2$, subestimando sistemáticamente la varianza poblacional. Dividir entre los $n-1$ grados de libertad corrige este sesgo exactamente.
\end{solproblema}

% Evaluar
\begin{solproblema}[prob:293fd20]
	\begin{enumerate}
		\item \textbf{Inexistencia de momentos:} la esperanza de $X$ con densidad $f(x) = \frac{1}{\pi(1+x^2)}$ está dada por $E(X) = \frac{1}{\pi}\int_{-\infty}^{\infty} \frac{x}{1+x^2}dx$. Para que la integral converja en el sentido de Lebesgue, las partes positiva y negativa deben ser absolutamente finitas por separado:
		\begin{align}
			\int_0^{\infty} \frac{x}{1+x^2} dx = \lim_{b \to \infty} \frac{1}{2}\ln(1+b^2) = \infty.
		\end{align}
		Como la integral sobre $[0,\infty)$ diverge (y análogamente sobre $(-\infty,0]$), la integral impropia total es una indeterminación $\infty - \infty$. Por tanto, $E(X)$ (y con mayor razón $\sigma^2$) \textbf{no existe}.

		\item \textbf{Distribución de la media muestral:} la función característica de un Cauchy estándar es $\varphi_X(t) = e^{-|t|}$. Para $\bar{X}_n = \frac{1}{n}\sum X_k$:
		\begin{align}
			\varphi_{\bar{X}_n}(t) = \prod_{k=1}^n \varphi_X\left(\frac{t}{n}\right) = \exp\left(-n\frac{|t|}{n}\right) = e^{-|t|}.
		\end{align}
		Esta función característica es idéntica a la de una sola observación, por lo que $\bar{X}_n \sim \text{Cauchy}(0,1)$ para cualquier $n$: \textbf{no ocurre ninguna convergencia hacia la normal} ni reducción de varianza. Un analista que asuma ciegamente el TLC (por ejemplo, aplicando la "regla del $n\ge30$" sin verificar que $\sigma^2<\infty$) sobre datos con colas extremadamente pesadas —comunes en rendimientos financieros o latencias de red— construiría intervalos de confianza y pruebas de hipótesis completamente inválidos, ya que promediar más datos no reduce el error de estimación en absoluto.
	\end{enumerate}
\end{solproblema}

% Crear
\begin{solproblema}[prob:b7567ec]
	\textbf{Ejemplo:} el diámetro de los tornillos producidos por una máquina CNC sigue una distribución con media poblacional $\mu = 10$ mm y desviación estándar $\sigma = 0.5$ mm. Se toma una muestra aleatoria de $n = 64$ tornillos, y se desea calcular $P(9.9 < \bar{X} < 10.1)$.

	El error estándar de la media es $\sigma_{\bar{X}} = \sigma/\sqrt{n} = 0.5/\sqrt{64} = 0.0625$ mm. Estandarizando los límites:
	\begin{align}
		z_1 = \frac{9.9-10}{0.0625} = -1.6, \qquad z_2 = \frac{10.1-10}{0.0625} = 1.6.
	\end{align}
	Por el TLC, $\bar{X}$ se aproxima a una normal, así que:
	\begin{align}
		P(9.9 < \bar{X} < 10.1) \approx P(-1.6 < Z < 1.6) = 2\Phi(1.6) - 1 = 2(0.9452) - 1 = 0.8904 \text{ (89.04\%).}
	\end{align}
\end{solproblema}
